% Copyright (c) 2026 Eugene Joseph M. Ragasa.
% SPDX-License-Identifier: Apache-2.0
% Author: Eugene Joseph M. Ragasa
% Formal proof completed: 2026-10-20
\documentclass[11pt]{article}
\usepackage[T1]{fontenc}\usepackage{amsmath,amssymb}\usepackage[margin=1in]{geometry}\usepackage[hidelinks]{hyperref}
\title{Proof Dossier: Rank-Two Frobenius Coefficient Distance}\author{Eugene Joseph M. Ragasa}\date{Formal proof completed: 2026-10-20}
\begin{document}\maketitle
\section{Identity and status}
ID \texttt{BG-FOUNDATION-THM-011}; status \texttt{proved}. Lean declaration:
\par\noindent{\tiny\path{ProjectKoios.BandGauge.Foundations.}}\\
{\tiny\path{rankTwoHermitianMatrix_frobeniusSq_sub}}.
\section{Source and contribution boundary}
Source-neutral supporting algebra; relation \texttt{background-only} to standard finite
Hermitian matrix theory through pinned Mathlib. No publication-source or novelty claim.
\section{Statement and proof}
For coefficient tuples $(t,x,y,z)$ and $(T,X,Y,Z)$, the unnormalized entrywise squared
Frobenius distance is
\[
2\bigl((t-T)^2+(x-X)^2+(y-Y)^2+(z-Z)^2\bigr).
\]
The proof expands all four complex entries and collects real squares.
\section{Boundary}
This is the unnormalized Frobenius quantity. It does not introduce the paper's
per-retained-state factor $1/2$. Kernel axioms are \texttt{propext}, \texttt{Classical.choice}, and \texttt{Quot.sound}; no admissions. Checked with Lean \texttt{4.35.0} and Mathlib revision \texttt{56a005855ea17e0b0958f73ec801ae140b0a1679}. AI-assisted and kernel-checked. No independent expert review is recorded.
\end{document}
